\section{格}

\begin{definition}{作为偏序集的格}{}
	设集合$E$配备偏序$\leq$.若对任意$x,y\in E$有$\uparrow_{E}^{\leq}\left(\left\{x,y\right\}\right)\neq\emptyset\wedge\downarrow_{E}^{\leq}\left(\left\{x,y\right\}\right)\neq\emptyset$,则称$E$为格,
\end{definition}

\begin{proposition}{}{}
	如果$\left(E_{i}\right)_{i\in I}$是一族格,则配备逐点序后$\prod\limits_{i\in I}E_{i}$是格.
\end{proposition}

\begin{definition}{幂集格}{}
	设$A$是集合,则$\mathcal{P}\left(A\right)$配备包含偏序后是格,称为$A$的幂集格.
\end{definition}

\begin{example}{}{}
	\begin{enumerate}
		\item 给$\Z_{\geq 1}$配备整除偏序,则对任意$x,y\in\Z_{\geq 1}$,$\sup\left(x,y\right)=\lcm\left(x,y\right),\inf\left(x,y\right)=\gcd\left(x,y\right)$.
		\item 设$G$是群,则$G$的所有子群组成的集合配备包含偏序后是格,称为$G$的子群格.
		\item 设$E$是集合,则$E$上的所有拓扑组成的集合配备包含偏序后是格,称为$A$的拓扑格.
		\item 设$I$是集合,则$\Hom_{\mathsf{Set}}\left(I,\R\right)$配备逐点序后是格,称为实值函数格.
	\end{enumerate}
\end{example}

\begin{remark}{}{}
	格一定是上定向集和下定向集,但既是上定向集又是下定向集的预序集不一定是格.
	
	例如,给$\left\{\ev_{p}\in\End_{\mathsf{Ring}}\left(\R\right)\middle|p\in\R\left[X\right]\right\}$配备逐点序,则该集合是上定向集和下定向集,但不是格.
\end{remark}










































